\documentclass[12pt]{article}
\usepackage{amsmath, amssymb, geometry}
\geometry{a4paper, margin=1in}
\setcounter{MaxMatrixCols}{100}

\begin{document}

\title{Phase Space Debug Report: \texttt{$5^2_1$} (SnapPy: \texttt{L5a1})}
\author{Automated Pipeline}
\date{\today}
\maketitle

\section{Gluing Equations (SnapPy)}
Manifold: \texttt{L5a1},\ $n = 4$ tetrahedra,\ $r = 2$ cusps.
Columns ordered as $Z_1, Z_1', Z_1'',\, Z_2, Z_2', Z_2'',\, Z_3, Z_3', Z_3'',\, Z_4, Z_4', Z_4''$.

\textbf{Tetrahedra equations ($C$ matrix, shape $4 \times 12$):}
\[ C = \begin{bmatrix}
1 & 0 & 0 & 1 & 0 & 1 & 1 & 0 & 1 & 1 & 0 & 0 \\
0 & 1 & 0 & 0 & 0 & 1 & 0 & 0 & 1 & 0 & 1 & 0 \\
0 & 1 & 1 & 0 & 1 & 0 & 0 & 1 & 0 & 0 & 1 & 1 \\
1 & 0 & 1 & 1 & 1 & 0 & 1 & 1 & 0 & 1 & 0 & 1
\end{bmatrix} \]

\textbf{Cusp equations (row order: $\mu_0,\, \lambda_0,\, \mu_1,\, \lambda_1$):}
\[ C_{\mathrm{cusp}} = \begin{bmatrix}
 0 &  0 &  0 &  0 & -1 &  0 &  0 &  0 &  1 &  0 &  0 & -1 \\
-1 &  1 &  1 &  0 & -2 &  0 & -1 &  1 &  1 &  0 &  0 & -2 \\
 0 &  0 &  0 &  0 &  0 &  0 &  0 & -1 &  0 &  1 &  0 &  0 \\
 0 &  1 &  0 &  0 &  0 & -1 &  0 &  0 &  1 &  0 & -1 &  0
\end{bmatrix} \]

\section{Geometric Basis}
The pipeline finds $n - r = 2$ independent internal edges.
Both are classified as \textbf{hard} edges (no easy edges were found for this triangulation).

\begin{itemize}
    \item \textbf{Hard Edge 1:}\ Vector $= [1,\, 0,\, 0,\, 1,\, 0,\, 1,\, 1,\, 0,\, 1,\, 1,\, 0,\, 0]$
    \item \textbf{Hard Edge 2:}\ Vector $= [0,\, 1,\, 0,\, 0,\, 0,\, 1,\, 0,\, 0,\, 1,\, 0,\, 1,\, 0]$
\end{itemize}

Since both internal edges are hard, there are \textbf{two fugacity variables} $\eta_1, \eta_2$ in the refined index.

\section{Neumann-Zagier Matrix}
Row structure of $g_{NZ} \in \mathrm{Sp}(8,\mathbb{Q})$
($2n \times 2n$ with $n=4$, columns in block order $Z_1, Z_2, Z_3, Z_4, Z_1'', Z_2'', Z_3'', Z_4''$):
\[
\text{rows } 0,1: \mu_0, \mu_1 \quad
\text{rows } 2,3: Q_1, Q_2 \quad
\text{rows } 4,5: \tfrac{\lambda_0}{2}, \tfrac{\lambda_1}{2} \quad
\text{rows } 6,7: \Gamma_1, \Gamma_2
\]

\textbf{$g_{NZ}$ Matrix:}
\[ g_{NZ} = \begin{bmatrix}
 0 &  1 &  0 &  0 &  0 &  1 &  1 & -1 \\
 0 &  0 &  1 &  1 &  0 &  0 &  1 &  0 \\
 1 &  1 &  1 &  1 &  0 &  1 &  1 &  0 \\
-1 &  0 &  0 & -1 & -1 &  1 &  1 & -1 \\
-1 &  1 & -1 &  0 &  0 &  1 &  0 & -1 \\
-\frac{1}{2} &  0 &  0 &  \frac{1}{2} & -\frac{1}{2} & -\frac{1}{2} &  \frac{1}{2} &  \frac{1}{2} \\
 0 & -1 &  0 & -1 &  0 &  0 &  0 &  0 \\
 1 &  0 &  \frac{1}{2} &  \frac{1}{2} &  0 &  0 &  \frac{1}{2} &  0
\end{bmatrix} \]

\textbf{Affine Shift $\nu$:}
\[
\nu_X = \begin{bmatrix} -1 \\ -1 \\ -2 \\ 0 \end{bmatrix}, \qquad
\nu_P = \begin{bmatrix} 0 \\ 0 \\ 0 \\ 0 \end{bmatrix}
\]

\section{Inversion and Local Charges}
\textbf{Inverse Matrix $g_{NZ}^{-1}$ (via symplectic identity $g^{-1} = [[D^T, -B^T],[-C^T, A^T]]$):}
\[ g_{NZ}^{-1} = \begin{bmatrix}
 0 & -\frac{1}{2} &  0 &  0 &  0 &  0 &  0 &  1 \\
 1 & -\frac{1}{2} &  0 &  0 & -1 &  0 & -1 & -1 \\
 0 &  \frac{1}{2} &  0 &  \frac{1}{2} & -1 & -1 & -1 & -1 \\
-1 &  \frac{1}{2} &  0 &  0 &  1 &  0 &  0 &  1 \\
 1 &  \frac{1}{2} &  0 & -1 &  0 &  0 &  1 & -1 \\
-1 &  0 &  1 &  0 &  1 &  0 &  1 &  0 \\
 1 &  0 &  0 & -\frac{1}{2} &  0 &  1 &  1 &  0 \\
 0 & -\frac{1}{2} &  1 & -\frac{1}{2} &  0 &  1 &  1 & -1
\end{bmatrix} \]

The master charge vector is
\[
\kappa = \bigl[m_1,\; m_2,\; 0,\; 0,\; e_1,\; e_2,\; e_{\mathrm{int},1},\; e_{\mathrm{int},2}\bigr]^T,
\]
where $m_1, m_2$ are cusp meridian charges, $e_1, e_2$ are cusp longitude$/2$ charges,
and $e_{\mathrm{int},1}, e_{\mathrm{int},2} \in \tfrac{1}{2}\mathbb{Z}$ are summed over.

The local tetrahedron charges $\Delta = g_{NZ}^{-1}\,\kappa$:
\begin{align*}
\Delta_1 &: \quad
  \tilde{m}_1 = -\tfrac{m_2}{2} + e_{\mathrm{int},2},
  \qquad
  \tilde{e}_1 = m_1 + \tfrac{m_2}{2} + e_{\mathrm{int},1} - e_{\mathrm{int},2} \\[4pt]
\Delta_2 &: \quad
  \tilde{m}_2 = m_1 - \tfrac{m_2}{2} - e_1 - e_{\mathrm{int},1} - e_{\mathrm{int},2},
  \qquad
  \tilde{e}_2 = -m_1 + e_1 + e_{\mathrm{int},1} \\[4pt]
\Delta_3 &: \quad
  \tilde{m}_3 = \tfrac{m_2}{2} - e_1 - e_2 - e_{\mathrm{int},1} - e_{\mathrm{int},2},
  \qquad
  \tilde{e}_3 = m_1 + e_2 + e_{\mathrm{int},1} \\[4pt]
\Delta_4 &: \quad
  \tilde{m}_4 = -m_1 + \tfrac{m_2}{2} + e_1 + e_{\mathrm{int},2},
  \qquad
  \tilde{e}_4 = -\tfrac{m_2}{2} + e_2 + e_{\mathrm{int},1} - e_{\mathrm{int},2}
\end{align*}

The topological framing phase power
$\phi = \sum_k \tilde{m}_k \cdot (\nu_P)_k - \sum_k \tilde{e}_k \cdot (\nu_X)_k$.
Since $\nu_P = 0$ this reduces to:
\[
\phi
= -\bigl(\tilde{e}_1(-1) + \tilde{e}_2(-1) + \tilde{e}_3(-2) + \tilde{e}_4 \cdot 0\bigr)
= \tilde{e}_1 + \tilde{e}_2 + 2\tilde{e}_3.
\]
Substituting the local charges:
\[
\phi = 2m_1 + \tfrac{m_2}{2} + e_1 + 2e_2 + 4e_{\mathrm{int},1} - e_{\mathrm{int},2}.
\]

\section{Final Summation Formula}
The refined index is:
\[
I_{5^2_1}(m_1, m_2, e_1, e_2)
= \sum_{e_{\mathrm{int}} \in \frac{1}{2}\mathbb{Z}^2}
   \Bigl(\prod_{j=1}^{2} \eta_j^{2e_{\mathrm{int},j}}\Bigr)
   \cdot (-q^{1/2})^{\phi}
   \cdot \prod_{i=1}^{4} I_{\Delta}(\tilde{m}_i,\, \tilde{e}_i).
\]

\textbf{Evaluated output at $(m_1, m_2, e_1, e_2) = (0,0,0,0)$ (truncated at $q^4$):}
\begin{align*}
I &= 1 \\
&\quad - 4q + q\,\eta_2^{2} \\
&\quad + q^2\!\left( -2 + \eta_1^{-1} + \eta_1 + \eta_2^{4} \right) \\
&\quad + q^3\!\left( 12 - 4\eta_2^{2} + \eta_2^{-2} + \eta_1^{-1}\eta_2^{-2} + \eta_1\eta_2^{-2}
                    + 2\eta_1^{-1} + 2\eta_1 + \eta_2^{6} \right) \\
&\quad + q^4\!\left( 17 - 8\eta_2^{2} + 4\eta_2^{-2} - \eta_1^{-1} + 2\eta_1^{-1}\eta_2^{-2}
                    - \eta_1 + 2\eta_1\eta_2^{-2} + \eta_2^{4} + \eta_2^{8} \right)
\end{align*}

\noindent\textbf{Raw key convention} (as stored in \texttt{RefinedIndexResult}):
a monomial $q^k\,\eta_1^a\,\eta_2^b$ is stored under key $(2k,\; 2a,\; 2b)$.

\bigskip
\noindent\textbf{Remark on triangulation difference:}
The reference document \texttt{5\_2\_1\_debug.tex} was generated from a different SnapPy
triangulation of $5^2_1$ (the pipeline found one easy and one hard edge there,
yielding a $1 \times 1$ fugacity sector).
The current pipeline's triangulation for \texttt{L5a1} has no easy edges,
so the refined index naturally carries two fugacity variables $\eta_1, \eta_2$.

\end{document}
